\chapter{Design and Architecture}
\label{ch:design}

\section{Introduction}

This chapter describes how the platform is put together. It begins with the
overall architecture and the principles that guided it, then explains the design
of each major subsystem: the grounded chat, the agent, and the knowledge graph.
It then presents the data model that ties everything together and the way the
system is deployed. The implementation detail that realises this design is the
subject of Chapter~5.

\section{Architectural Overview}

Synapse is organised into tiers, as shown in Figure~\ref{fig:architecture}. The
web application runs in the browser and talks to a Cloudflare Worker, which hosts
the application at the edge and forwards requests to the backend. The backend is a
FastAPI service running on a GPU-equipped server, and it is where the real work
happens: an API layer receives requests, a pool of workers processes documents
through the pipeline, a retrieval and reasoning component answers questions, and a
knowledge-graph builder constructs the graph. These components share two stores,
a Supabase database with vector search for text and embeddings and a Cloudflare
R2 bucket for files and artefacts, and they call external model services for
language and vision tasks.

\begin{figure}[H]
  \centering
  \includegraphics[width=\textwidth]{fig_architecture.png}
  \caption{High-level architecture of Synapse. The browser talks to a Cloudflare
  Worker that hosts the application and proxies to a FastAPI service on a
  GPU server. The backend separates the API layer, the document workers, the
  retrieval and reasoning component and the knowledge-graph builder, all sharing a
  Supabase database and an R2 object store and calling external model services.
  This separation lets each concern scale and evolve on its own.}
  \label{fig:architecture}
\end{figure}

\section{Design Principles}

A small number of principles shaped the design. The first is \emph{grounding}:
every answer must be traceable to a source passage, and the system must be able to
say it does not know. The second is \emph{separation of concerns}: ingestion,
retrieval, artefact generation and graph construction are independent components
with clear interfaces, which kept each one simple and allowed them to be built and
tested in turn. The third is \emph{resumable processing}: documents move through
the pipeline in stages that can be retried, so a failure never discards completed
work. The fourth is \emph{tenant isolation}: every piece of data belongs to an
organisation, and access is enforced at the database rather than trusted to the
application. These principles recur throughout the subsystems described below.

\section{The Grounded Chat Subsystem}

The grounded chat answers questions using only the user's selected libraries.
Its design, shown in Figure~\ref{fig:chatflow}, follows the corrective and
multi-agent ideas from Chapter~2. A question is first analysed and, where useful,
rewritten and expanded so that retrieval has the best chance of finding the right
passages. Retrieval is hybrid: a vector search finds passages by meaning and a
keyword search finds them by exact terms, and the two ranked lists are fused. The
fused candidates are re-ranked by a cross-encoder, and the top passages are graded
for relevance and sufficiency. Only if the evidence is sufficient does the system
compose an answer, attaching a citation to each claim; otherwise it abstains and
says the library cannot support an answer.

\begin{figure}[H]
  \centering
  \includegraphics[width=0.62\textwidth]{fig_chat_flow.png}
  \caption{Design of the grounded chat. The question is analysed, evidence is
  retrieved by both meaning and keyword and then re-ranked, and the evidence is
  graded before any answer is written. A sufficient set of passages leads to a
  cited answer; an insufficient set leads the system to abstain rather than guess.
  This grade-then-answer order is what keeps the replies honest.}
  \label{fig:chatflow}
\end{figure}

This design directly serves the honesty requirement. Because the decision to
answer is made after the evidence has been graded, the system can decline cleanly,
and because each claim is tied to a passage, a reader can verify the answer
against the source.

\section{The Agent Subsystem}

The agent turns a goal and a set of libraries into an artefact such as a chart, a
document or an image. Its design is shown in Figure~\ref{fig:agentflow}. The agent
first plans the work, deciding what the user wants and what data is required. If
the request is genuinely ambiguous it asks a single clarifying question;
otherwise it proceeds and records the assumptions it has made. It then acquires
the data, either by re-parsing a structured file exactly or by retrieving from the
library, builds a specification of the artefact, validates and if necessary
repairs it, renders it both interactively and as a downloadable file, and finally
explains the result.

\begin{figure}[H]
  \centering
  \includegraphics[width=\textwidth]{fig_agent_flow.png}
  \caption{Design of the agent. A goal is planned into data needs and artefacts;
  the agent clarifies only when the request is ambiguous, acquires the data,
  builds and validates a specification, and renders the artefact for viewing and
  download. The data values are bound to a chart in code rather than written by
  the model, so the figures always match the source document.}
  \label{fig:agentflow}
\end{figure}

The key design decision here is the separation between describing an artefact and
filling it with data. The model proposes the shape of a chart, but the numbers are
bound to it from the source file in code. This keeps the artefacts faithful, which
is essential if a user is to rely on them.

\section{The Knowledge-Graph Subsystem}

The knowledge-graph subsystem reveals how the ideas in a library relate. Its
construction is shown in Figure~\ref{fig:kgbuild}. Starting from the processed
library, it extracts entities and relationships from each chunk using a language
model, resolves duplicate entities so that the same thing mentioned in different
ways becomes a single node, builds the graph by assembling the nodes and edges and
weighting them, and serves the result as an interactive graph in the web
application. Every node and edge keeps a link back to the passage it came from, so
the graph is not only a picture but also a way back into the evidence.

\begin{figure}[H]
  \centering
  \includegraphics[width=\textwidth]{fig_kg_build.png}
  \caption{Design of the knowledge-graph subsystem. Entities and relationships are
  extracted from each chunk, duplicate entities are resolved into single nodes, the
  graph is assembled and weighted, and it is served interactively. Provenance is
  kept throughout, so a user can move from a node or edge back to the source
  passage that supports it.}
  \label{fig:kgbuild}
\end{figure}

\section{Data Model}

The data model, shown in Figure~\ref{fig:datamodel}, ties the subsystems together.
At its centre is the organisation, which is the tenant boundary; users belong to an
organisation, and libraries are owned by it. A library contains documents, and each
document is split into chunks that carry both their text and their embedding, which
is what retrieval searches. Conversations and their messages, and the nodes and
edges of the knowledge graph, are likewise attached to the organisation. Every row
is scoped to an organisation, and access is enforced by row-level security in the
database, which is how tenant isolation is guaranteed.

\begin{figure}[H]
  \centering
  \includegraphics[width=\textwidth]{fig_data_model.png}
  \caption{Simplified core data model. The organisation is the tenant boundary;
  users, libraries, documents, chunks, conversations and graph elements all belong
  to it. Chunks carry the embeddings that retrieval searches, and the graph
  elements are grounded in the same chunks, so the data model supports both
  answering and mapping from one consistent structure.}
  \label{fig:datamodel}
\end{figure}

\section{Deployment Architecture}

Figure~\ref{fig:deployment} shows how the parts are deployed. The web application
is served from a Cloudflare Worker, with files and artefacts in an R2 bucket and
live updates delivered over a realtime channel. The heavy processing runs on a
single GPU-equipped virtual machine that hosts the FastAPI service, the pipeline
workers and the local embedding and reranking models. The database, with its
vector search, is a managed Supabase instance, and the language and vision models
are reached over the network. This arrangement keeps the user-facing parts fast
and close to the user while concentrating the graphics-intensive work on one
server.

\begin{figure}[H]
  \centering
  \includegraphics[width=\textwidth]{fig_deployment.png}
  \caption{Deployment view. The application and its object store run on Cloudflare
  at the edge, while document processing, retrieval and the local models run on a
  single GPU virtual machine. The managed database provides vector search, and the
  external model services are called over the network. The split keeps the
  interface responsive and the processing in one place.}
  \label{fig:deployment}
\end{figure}

\section{Summary}

This chapter presented the architecture of Synapse and the principles of
grounding, separation of concerns, resumable processing and tenant isolation that
shaped it. It described the design of the grounded chat, which grades evidence
before answering and abstains when it is insufficient; the agent, which builds
faithful artefacts by binding data in code; and the knowledge graph, which keeps
provenance back to the source. It then presented the data model and the deployment
arrangement. The next chapter explains how this design was implemented.
